\section{The original local Hecke operators, projected lattice and spectral record}
\label{sec:local-hecke-projection}

\status{HLP1--HLP32 reconstruct the coefficient-linear Hecke lift,
its integral local standard facts, and the actual localized lattice,
projection, faithful action and complete \texttt{nonempty\_\allowbreak spec}
assembly. Every statement retains its printed hypotheses. In particular,
the complex Hecke standard facts, the original rational and integral
\texttt{Mt\allowbreak Input}, and the compatible arithmetic coefficient and character
suppliers are not inferred from the existence of an abstract commuting
algebra. This chapter supplies the finite algebraic part of the local
construction; it does not yet assert the unconditional producer
\texttt{nd\allowbreak Spec\_\allowbreak exists}, \texttt{Aut\allowbreak Spec} or \texttt{NBuilt}.}

We retain the original \(D:\texttt{D7}\), its family \(M\), normalized
choice \(C\), and fixed complex embedding
\(\epsilon:L=D.\mathrm{what}_2.\mathrm{LL}\hookrightarrow\mathbb C\).
For an eligible pair \((j,w)\), write
\[
 U=D.\mathrm{what}_2.\mathrm{MtL}(\mu^{\rm tw}_w,j),\quad
 j_M:U\longrightarrow D.\mathrm{what}_2.\mathrm{Mt}(\mu^{\rm tw}_w,j),
 \quad e=M.\mathrm{real}\circ j_M.
\]
The map \(j_M\) is semilinear for \(\epsilon\), carries the original
basis to the original complex basis, and is injective. The range of
\(M.\mathrm{real}\) is exactly its specified \(\operatorname{TypeSp}\).
After a specified coefficient extension \(X\), let \(O=O_{j,w}\) be its
actual complete DVR, \(K=\operatorname{Frac}(O)\),
\(e_c:L\hookrightarrow K\) its specified map, and
\(V=K\otimes_{L,e_c}U\). No coefficient embedding is replaced.

% BEGIN HLP1
\paragraph{HLP1: the full ambient generator range and fixed lattice.}
\label{claim:HLP1}
The predicate \(\operatorname{GenAt}(D,\ell,C.Q(w),v,k)\) means that
\(v\) is split, is not above \(\ell\), belongs to neither the original
exceptional set nor \(C.Q(w)\), and \(1\le k\le n_V\).
This is definitionally \(\operatorname{GenAt0}\) on \(\texttt{base\allowbreak C}\).
The upper index is \(n_V\), rather than the row index bound \(r\) in
LTD6. Throughout this chapter the original \(\texttt{Mt\allowbreak Input.rat}\) and
\(\texttt{Mt\allowbreak Input.int}\) are required for this entire ambient range.

Let \(G\) denote the original finite set
\(D.\mathrm{what}_2.\mathrm{MtLat.gen}(\mu^{\rm tw}_w,j)\), and set
\[
       S=j_M^{-1}(G),\qquad
       L_0=\operatorname{span}_O\{i_V(s):s\in S\}\subset V,
       \quad i_V(u)=1\otimes u.                     \tag{HLP.1}
\]
The supplied row datum states \(G\subset\operatorname{range}(j_M)\)
and \(\operatorname{span}_{\mathbb C}G=\mathrm{Mt}\). Consequently
\(S\) is exactly the source's \texttt{lat\allowbreak L}, including its given
finite generating elements. Later projected lattices will be spans
of these elements, not arbitrary lattices in the same vector space.
The actual \(Q=C.Q(w)\), bad sets and uniformizers occur in every
operator formula below. This checks the object dictionary used in
\texttt{Hecke\allowbreak Lift0}, \texttt{NPart\allowbreak CNd} and \texttt{Nd\allowbreak Spec}.
% END HLP1

% BEGIN HLP2
\paragraph{HLP2: independence of the uniformizer and the permitted split-place witness.}
\label{claim:HLP2}
Fix an original split place \(s\), and two uniformizers \(\varpi,\varpi'\)
at its real place. Their ratio \(u=\varpi/\varpi'\) and its inverse
have valuation zero. The diagonal matrix with its first \(k\) entries
\(u\) and its other entries \(1\) therefore belongs to
\(\operatorname{GL}_{n_V}(O_{s.v})\); its inverse has the same property.
Multiplying the diagonal Hecke representative for \(\varpi'\) by this
matrix gives the representative for \(\varpi\). The two double cosets,
and hence the source's sets \(\texttt{t\allowbreak Set}\), agree. Their Hecke sums
agree as functions because the double-coset operator is defined using
the same coset set. The normalizing factor \(q_s^{-k(n_V-1)}\) also
agrees.

If \(\operatorname{IsConjPl}(v,s.w)\) and
\(\operatorname{IsConjPl}(v,s'.w)\), the defining conjugate place is
unique, so \(s.w=s'.w\). We check that this also identifies the
embedding field of the two \texttt{Split\allowbreak Place} records. Their real
places agree by restriction. Write their embeddings into the same
completion as \(i,i':L\to L^+_{s.v}\). The valuation ring that either
induces on \(L\) is \((\mathcal O_L)_{s.w}\). Indeed the integers
have norm at most one, and those of norm less than one form exactly
the prime \(s.w\). Elements outside that prime have norm one, so
its localization is contained in the induced valuation ring.
Conversely, if \(x\) is outside that localization, its inverse
belongs to the maximal ideal of this local DVR. Write that inverse
as a fraction with numerator in \(s.w\) and denominator outside
\(s.w\); its norm is less than one. Thus \(x\) cannot belong to
the induced valuation ring. Both embeddings therefore satisfy
\(|i(x)|\le1\) if and only if \(|i'(x)|\le1\), for every \(x\in L\).

If \(\delta=i(x)-i'(x)\ne0\), density of \(L^+\) in its completion
gives \(a,b\in L^+\) with
\(|i(x)-a|<|\delta|\) and \(|\delta-b|<|\delta|\).
The ultrametric law gives \(|b|=|\delta|>0\). For
\(z=(x-a)/b\), it follows that \(|i(z)|<1\) and
\(|i'(z)|=1\). Applying the equality of valuation rings to
\(z^{-1}\) is a contradiction. Hence \(i=i'\), and the records
are equal. This supplies the embedding-uniqueness argument in
\texttt{split\allowbreak Place\_\allowbreak eq\_\allowbreak of\_\allowbreak w\_\allowbreak eq}, rather than treating the
embedding as an automatically canonical record field.

The uniformizer argument now applies to \(s,s'\) and proves the
source's \texttt{t\allowbreak N\_\allowbreak eq\_\allowbreak of\_\allowbreak conj}. In particular, a lift constructed
with one permitted witness satisfies the equation for every
permitted witness. This does not identify the operators at two
unrelated places.
% END HLP2

% BEGIN HLP3
\paragraph{HLP3: the complete coefficient-linear lift, before all split witnesses.}
\label{claim:HLP3}
Assume the original \texttt{Hecke\allowbreak Std\allowbreak Plain}, \texttt{Gen\allowbreak Hyp} and
\texttt{Mt\allowbreak Input}, with the same \(D,M,C\) and tuple. For fixed eligible
\(j,w,v,k\), we construct one \(L\)-linear endomorphism \(T\) of \(U\)
such that
\[
 e(Tu)=\operatorname{tN}(M.J,s,\varpi,k)e(u)
   \quad\hbox{for every permitted }s,\varpi\hbox{ and every }u. \tag{HLP.2}
\]
If there is no such pair \((s,\varpi)\), take \(T=0\); the universal
condition is empty. Otherwise fix one pair. The exact rationality
clause \texttt{Mt\allowbreak Input.rat} puts the right side of (HLP.2) in
\(e(U)\). Injectivity of \(e\) gives a unique preimage for every \(u\).
The complex-linearity clauses of \texttt{Hecke\allowbreak Std\allowbreak Plain} apply on this
\(\operatorname{TypeSp}\), since \texttt{Gen\allowbreak Hyp} excludes the same
place from its bad set. For \(a\in L\), they show that the preimages
of the sum and of \(\epsilon(a)e(u)\) are respectively the sum and
\(a\) times the preimage. Injectivity proves additivity and
\(L\)-linearity. HLP2 gives (HLP.2) for all the other witnesses.

This proves the full quantifier order in \texttt{Hecke\allowbreak Lift0}: the
operator is chosen before \(s,\varpi\), and its domain is the original
\(U\). It does not manufacture \texttt{Mt\allowbreak Input.rat}; that assertion
about the actual form realization is an antecedent here.
% END HLP3

% BEGIN HLP4
\paragraph{HLP4: the chosen operators commute at every index and base change literally.}
\label{claim:HLP4}
Define \(h_L(j,w,v,k)\) to be the chosen lift of HLP3 when both
\(1\le j\le m\) and \(\operatorname{GenAt}(v,k)\) hold, and zero
otherwise. This is the literal \texttt{h\allowbreak LC}. For two eligible
generators choose permitted split-place and uniformizer witnesses,
as supplied by the original split-place existence assertion.
Apply \(e\) to the two composites. Twice using (HLP.2) turns them
into the two original complex Hecke composites on \(e(U)\).
\texttt{Hecke\allowbreak Std\allowbreak Plain}, with the common bad set of
\texttt{Gen\allowbreak Hyp}, says they agree. Injectivity of \(e\) proves
\([h_L(v,k),h_L(v',k')]=0\). If an index is ineligible, one operator
is zero and the same conclusion holds. This proves the all-index
statement \texttt{h\allowbreak LC\_\allowbreak commute}, including its guarded zero branches.

Put \(h_K=1\otimes h_L\) on the actual \(V=K\otimes_LU\). Then
\[
       h_K i_V=i_Vh_L.                              \tag{HLP.3}
\]
Evaluation on pure tensors proves both (HLP.3) and commutativity
after base change. The \(K\)-span of \(i_V(U)\) is \(V\), so these
formulas uniquely determine \(h_K\). They are exactly
\texttt{h\allowbreak KC\_\allowbreak def} and the operator fields of \texttt{base\allowbreak C}.
% END HLP4

% BEGIN HLP5
\paragraph{HLP5: finite dimension and the actual characteristic polynomials.}
\label{claim:HLP5}
The given basis of \(U\) maps under \(j_M\) to a basis of the finite
complex space \(\mathrm{Mt}\). Thus it is finite, and both \(U\) and
\(V\) are finite-dimensional. The complexification \(T_{\mathbb C}\)
of an \(L\)-endomorphism \(T\) is defined by its matrix in this basis.
It satisfies \(T_{\mathbb C}j_M=j_MT\), and its characteristic
polynomial is \(\epsilon(\operatorname{charpoly}T)\), coefficient by
coefficient. Similarly \(\operatorname{charpoly}(1\otimes T)\) is
\(e_c(\operatorname{charpoly}T)\). These identities follow directly
by expanding the determinant, so they retain the chosen embeddings.

For \(T=h_L(v,k)\), (HLP.2) on a basis and complex linearity on
\(\operatorname{TypeSp}\) give
\[
 M.\mathrm{real}\circ T_{\mathbb C}
   =\operatorname{tN}(M.J,s,\varpi,k)\circ M.\mathrm{real}. \tag{HLP.4}
\]
A complex root of the characteristic polynomial has a nonzero
eigenvector, since the matrix is then singular; its image under
\(M.\mathrm{real}\) remains nonzero. A nonzero generalized
\(\mu\)-eigenspace of \(h_K\) also gives a nonzero eigenvector: from
\((h_K-\mu)^a x=0\), choose the smallest positive exponent and use
\((h_K-\mu)^{a-1}x\). Hence such a \(\mu\) is a root of the mapped
characteristic polynomial. When \(V=0\), the polynomial is \(1\),
no eigenvalue occurs, and these assertions are vacuous as required.
% END HLP5

% BEGIN HLP6
\paragraph{HLP6: prime-to-\(\ell\) complex denominators imply integrality in the actual DVR.}
\label{claim:HLP6}
Let \(p\in L[t]\) be monic, and suppose that for every complex root
\(\rho\) of \(\epsilon(p)\) there is an integer \(d_\rho\), prime to
\(\ell\), with \(d_\rho\rho\) integral over \(\mathbb Z\). Choose the
product \(c\) of these finitely many integers, using \(c=1\) when
there are no roots. Then \(\ell\nmid c\), and every \(c\rho\) is an
algebraic integer. If \(d=\deg p\), the monic polynomial
\[
             q(t)=c^d p(t/c)                         \tag{HLP.5}
\]
has coefficients in \(L\), and the coefficients of \(\epsilon(q)\)
are algebraic integers: each is an elementary symmetric polynomial
in the \(c\rho\), counted with multiplicity. The algebraic integers
are closed under sums and products (adjoin the finitely many
integral elements and use the determinant of their multiplication
matrices). Injectivity of \(\epsilon\) carries their monic integer
equations back to the coefficients of \(q\). Mapping by \(e_c\)
shows that these coefficients are integral over \(O\). They already
belong to \(K\), and a DVR is integrally closed: if a putative
integral element had negative valuation, its highest power in a
monic equation would have uniquely smallest valuation. Thus they
belong to \(O\).

If \(\mu\in K\) is a root of \(e_c(p)\), then \(c\mu\) is a root
of the monic polynomial \(e_c(q)\in O[t]\). The same argument gives
\(c\mu\in O\). The integer \(c\) is a unit in the supplied
\(\ell\)-adic DVR: it is already a unit in \(\mathbb Z_\ell\), and
the source algebra map preserves its inverse. Hence \(\mu\in O\).
No algebraic integrality at \(\ell\) was assumed for the individual
complex eigenvalues before the denominator was removed.
% END HLP6

% BEGIN HLP7
\paragraph{HLP7: the full original \texttt{Hecke\allowbreak Std0} on \texttt{base\allowbreak C}.}
\label{claim:HLP7}
Retain HLP3--HLP4's inputs and the original \texttt{Split\allowbreak K} supplied
by the chosen coefficient extension. For an eligible generator,
HLP5 sends every complex root of its rational characteristic
polynomial to a nonzero eigenfunction of the original complex
\(\operatorname{tN}\). Its place is not above \(\ell\): this follows
from \(\operatorname{GenAt}\) and the common real place in
\(\operatorname{IsConjPl}\). The last clause of
\texttt{Hecke\allowbreak Std\allowbreak Plain} consequently gives the prime-to-\(\ell\)
denominator in HLP6 for every such root. HLP6 and HLP5 now imply
that every \(K\)-eigenvalue with nonzero generalized eigenspace
belongs to the image of \(O\hookrightarrow K\). Together with
HLP4's commutation this is precisely both conjuncts of
\texttt{Hecke\allowbreak Std0}, for every \(1\le j\le m\). HLP5 gives
\texttt{Fin\allowbreak Dim\allowbreak V} for all \(j,w\), including the ineligible indices.

The splitting assertion in \texttt{Split\allowbreak K} is a separate field of
\texttt{base\allowbreak C}, not a consequence of eigenvalue integrality. MR23
constructs a finite extension splitting the possibly infinite
commuting family; when that construction is used it keeps the same
embedding \(e_c\). Thus the three inputs needed below---commutation,
splitting and integrality---have exactly their original provenance.
% END HLP7

% BEGIN HLP8
\paragraph{HLP8: integral Hecke combinations on the prescribed generators.}
\label{claim:HLP8}
Assume \(w\) is twisted and \(x\in S\). The original
\texttt{Mt\allowbreak Input.int}, applied to \(j_Mx\in G\), gives coefficients
\(c_g\in L\), each integral at \(w\), and the equality of complex
realizations
\[
 \operatorname{tN}e(x)
   =M.\mathrm{real}\!\left(\sum_{g\in G}\epsilon(c_g)g\right).
                                                               \tag{HLP.6}
\]
Only the indicated finite support need be retained. By (HLP.2) and
injectivity of \(M.\mathrm{real}\), the right-hand vector is
\(j_M(h_Lx)\). Since \(G\subset\operatorname{range}(j_M)\), each
supported \(g\) has one preimage \(y_g\in S\), and different \(g\)
have different preimages. Semilinearity and injectivity of \(j_M\)
therefore give the literal equality
\[
                 h_Lx=\sum_g c_g y_g.                \tag{HLP.7}
\]
This proves the complete original \texttt{Hecke\allowbreak Int\allowbreak Tw}, with its
finite subset of \texttt{lat\allowbreak L}, not a lattice chosen afterward.
The valuation compatibility of \(e_c\) converts each \(c_g\) to
an element of \(O\). Equation (HLP.3) then shows that every eligible
\(h_K\) preserves the original \(L_0\) of (HLP.1). The algebra they
generate over \(O\) preserves \(L_0\) as well, by closure under
addition, multiplication and scalar multiplication.
% END HLP8

% BEGIN HLP9
\paragraph{HLP9: every maximal ideal is one of the actual residual eigenspace classes.}
\label{claim:HLP9}
Fix an eligible \((j,w)\), and let
\(A=O[h_K(v,k):\operatorname{GenAt}(v,k)]\subset\operatorname{End}_K(V)\).
By HLP7 and LTD3 there are finitely many distinct integral characters
\(\psi_a:A\to O\), with nonzero simultaneous generalized spaces
\(W_a\), and \(V=\bigoplus_a W_a\). These are exactly the characters
with nonzero such space. For completeness, if \(\psi\) is another,
choose for each \(a\) an operator whose \(\psi\)- and
\(\psi_a\)-values differ. On \(W_a\) that operator minus its
\(\psi\)-value is invertible (a nonzero scalar plus a nilpotent).
No nonzero \(\psi\)-generalized vector can therefore have a
component in any \(W_a\).

Consider \(E:A\to O^s\), \(t\mapsto(\psi_a(t))_a\).
Every element of \(\ker E\) is nilpotent on each \(W_a\), and
hence nilpotent on their finite direct sum. Conversely a nilpotent
element maps to zero in every field of fractions. Thus \(\ker E\)
is the nilradical and belongs to every maximal ideal. Put
\(A'=E(A)\). It is an \(O\)-submodule of finite free \(O^s\), hence
finite, and it contains the diagonal scalars. The ring \(O^s\) is
finite as an \(A'\)-module, because its standard coordinate vectors
already generate it over the smaller scalar ring \(O\). It is
therefore integral over \(A'\).

We recall why a maximal ideal \(\mathfrak n\) of \(A'\) has a
maximal ideal above it in \(O^s\). Localize the finite injective
extension at \(A'\setminus\mathfrak n\). If its reduction modulo
\(\mathfrak n\) were zero, the determinant argument for a finite
module generated by \(x_1,\ldots,x_d\) with
\(x_i=\sum_j a_{ij}x_j\), \(a_{ij}\in\mathfrak n\), would give
\((1-b)O^s_{\mathfrak n}=0\) with \(b\in\mathfrak n\).
This contradicts the injected local ring. The reduction is a
nonzero finite-dimensional algebra over the residue field. A
maximal proper ideal of it, pulled back, gives an ideal over
\(\mathfrak n\); its residue field is a finite extension, so the
pulled-back ideal is maximal in the original integral extension.
The maximal ideals of \(O^s\) are exactly the coordinate ideals
with that coordinate in \(\mathfrak m_O\): its coordinate
idempotents force a field quotient to keep just one coordinate.
It follows that every maximal ideal of \(A\) is
\[
              \mathfrak m_a=\psi_a^{-1}(\mathfrak m_O). \tag{HLP.8}
\]
In particular there are finitely many, and each has a nonzero
supported character. If \(V=0\), \(A\) is the zero ring, has no
maximal ideals, and all these conclusions have the required empty
interpretation. We have not claimed that \(A\) itself is finite
over \(O\): its nilpotent part has not yet been bounded.
% END HLP9

% BEGIN HLP10
\paragraph{HLP10: the literal local indices, points and separating generators.}
\label{claim:HLP10}
The source's \(\texttt{Loc\allowbreak T}(B)\) consists of actual triples
\((j,w,\mathfrak m)\), with \(1\le j\le m\), \(w\) twisted,
and \(\mathfrak m\) a maximal ideal of that exact algebra \(A\).
For fixed \(j,w\), HLP9 makes this a finite set. Its map
\(\texttt{mx\allowbreak C}\) sends a matching triple to its own ideal and
an unmatched triple to zero. Equality of two matching ideals is
equality of the triples, so this map is injective on the matching
fiber, and every maximal ideal occurs there.

For such a triple, the literal set \(X_{\mathfrak m}=\texttt{XOf}\)
is the set of those \(\psi_a\) whose residual ideal is
\(\mathfrak m\). It is finite and nonempty by HLP9. The algebra
\(T_{\mathfrak m}=\texttt{TOf}\subset O^{X_{\mathfrak m}}\) is
 generated by the actual functions
\(h_T(v,k)(\psi)=\psi(h_K(v,k))\). Distinct points have distinct
characters. If all these generating values agreed, additivity,
multiplicativity and equality on \(O\) would imply agreement on
the algebra they generate, namely all of \(A\). Thus the displayed
generators separate the points. This proves the exact finite,
nonempty, maximal-ideal-index and separation fields, rather than
selecting an abstract set of eigenvalues with the same size.
% END HLP10

% BEGIN HLP11
\paragraph{HLP11: the original finite generator set spans before and after base change.}
\label{claim:HLP11}
Injectivity of \(j_M\) and finiteness of \(G\) make
\(S=j_M^{-1}(G)\) finite. It spans \(U\) over \(L\). Indeed, if
its span were proper, extend a basis of that span to one of \(U\)
and choose a nonzero \(L\)-linear functional vanishing on the
span. Extending this functional by \(\epsilon\) in the given
basis gives a nonzero complex functional on \(\mathrm{Mt}\)
vanishing on every \(j_Ms\). But these are all of \(G\), whose
complex span is the whole space. This contradiction proves
\(\operatorname{span}_LS=U\). Tensoring its spanning expressions
shows
\[
                   \operatorname{span}_K i_V(S)=V.    \tag{HLP.9}
\]
This verifies the hypothesis used by the source's
\texttt{Is\allowbreak Localized\allowbreak Of\allowbreak Span}; finite-dimensionality alone would
not justify (HLP.9) for an unrelated lattice.
% END HLP11

% BEGIN HLP12
\paragraph{HLP12: the actual projected module is finite and free over its DVR.}
\label{claim:HLP12}
For a local index \(\mathfrak m\), LTD3 and HLP9 give the original
direct sum
\[
 V=P_{\mathfrak m}\oplus Q_{\mathfrak m},\qquad
 P_{\mathfrak m}=\operatorname{NCov.part}(A,\mathfrak m),\quad
 Q_{\mathfrak m}=\operatorname{NCov.off}(A,\mathfrak m).
\]
Let \(p_{\mathfrak m}\) be projection along this exact complement.
Define
\[
 N_{\mathfrak m}
   =\operatorname{span}_O\{p_{\mathfrak m}i_V(s):s\in S\}
       \subset P_{\mathfrak m}.                      \tag{HLP.10}
\]
This is literally \texttt{Mt\allowbreak Sub}. It is finite because the set in
braces is finite, and torsion free because it is contained in a
\(K\)-vector space and \(O\hookrightarrow K\) is injective.
Here is the needed freeness argument without an additional
structure theorem. Choose \(K\)-coordinates on \(P_{\mathfrak m}\)
and multiply all generators by one nonzero common denominator.
The resulting isomorphic module is a submodule of \(O^d\).
Induct on \(d\): its image in the first coordinate is an ideal of
the DVR, hence is either zero or \(aO\). In the latter case choose
a vector with first coordinate \(a\). Subtracting a unique multiple
of this vector splits off a free rank-one summand, leaving a
submodule of \(O^{d-1}\). In the zero case use the induction
directly. This proves finite freeness and applies also to \(d=0\).
It establishes the actual \texttt{mt\allowbreak Sub\_\allowbreak finite} and
\texttt{mt\allowbreak Sub\_\allowbreak free} fields.
% END HLP12

% BEGIN HLP13
\paragraph{HLP13: localization is an isomorphism onto the actual part.}
\label{claim:HLP13}
Applying the surjective projection \(p_{\mathfrak m}\) to (HLP.9)
shows that \(N_{\mathfrak m}\) spans \(P_{\mathfrak m}\) over \(K\).
The canonical map
\[
 \eta:K\otimes_ON_{\mathfrak m}\longrightarrow P_{\mathfrak m},
                  \qquad a\otimes n\longmapsto an    \tag{HLP.11}
\]
is therefore surjective. It is injective by clearing denominators:
a finite sum of pure tensors can be written \(s^{-1}\otimes n\),
with \(0\ne s\in O\), and if its image is zero then \(n=0\)
in the actual submodule. The same denominator argument is exactly
the universal localization property for the inclusion
\(N_{\mathfrak m}\hookrightarrow P_{\mathfrak m}\): every
nonzero scalar from \(O\) acts invertibly after extending to \(K\).
Thus the source's \texttt{Mt\allowbreak Loc} holds, and \(\eta\) is its
specified \texttt{mt\allowbreak Equiv}. This argument requires no choice of
an eigenbasis and no semisimplicity.
% END HLP13

% BEGIN HLP14
\paragraph{HLP14: the literal comparison map, its exact range and its projection error.}
\label{claim:HLP14}
Define
\[
 \Psi=\bigl(P_{\mathfrak m}\hookrightarrow V\bigr)\circ\eta,
 \qquad \operatorname{pr}_{Mt}=\eta^{-1}p_{\mathfrak m}i_V.
                                                               \tag{HLP.12}
\]
These are the source's \texttt{Psi\allowbreak C} and \texttt{pr\allowbreak Mt\allowbreak C}; all scalar
structures are the original ones. HLP13 and injectivity of the
subtype inclusion prove that \(\Psi\) is injective and has range
exactly \(P_{\mathfrak m}\). Its defining expression gives
\[
       \Psi(\operatorname{pr}_{Mt}u)-i_Vu
         =p_{\mathfrak m}(i_Vu)-i_Vu\in Q_{\mathfrak m}. \tag{HLP.13}
\]
The last membership follows by decomposing \(i_Vu\) into its
\(P\)- and \(Q\)-components. This is the original error term
\texttt{Psi\_\allowbreak pr}; it is not an assertion that the projection equals
the identity. On \(s\in S\), (HLP.11) yields the exact integral
identity
\[
       \operatorname{pr}_{Mt}(s)=1\otimes p_{\mathfrak m}i_V(s).
                                                               \tag{HLP.14}
\]
% END HLP14

% BEGIN HLP15
\paragraph{HLP15: both required span assertions for the original projection.}
\label{claim:HLP15}
By (HLP.10) and (HLP.14), the \(O\)-span of
\(\operatorname{pr}_{Mt}(S)\) is precisely the image of
\(N_{\mathfrak m}\to K\otimes_ON_{\mathfrak m}\),
\(n\mapsto1\otimes n\). Both containments follow by applying this
map to a finite \(O\)-linear combination of the defining
generators. In particular every original integral generator maps
into the prescribed lattice.

The \(K\)-span of the entire range of \(\operatorname{pr}_{Mt}\)
is the full tensor space: HLP11 gives the span of \(i_V(U)\),
projection maps it onto \(P_{\mathfrak m}\), and \(\eta^{-1}\)
is an isomorphism. These are exactly \texttt{pr\allowbreak Mt\_\allowbreak lat} and
\texttt{pr\allowbreak Mt\_\allowbreak span}, together with the stronger integral span
identity proved in \texttt{NPart\allowbreak CNd\allowbreak Proj}. The proof keeps the
original rational domain and its original generating set.
% END HLP15

% BEGIN HLP16
\paragraph{HLP16: semisimplicity transfers from the original no-Jordan assertion.}
\label{claim:HLP16}
Assume the source's additional complex statement: on each actual
\(\operatorname{TypeSp}\),
\((\operatorname{tN}-\mu)^2f=0\) implies
\((\operatorname{tN}-\mu)f=0\), for every \(\mu\in\mathbb C\).
It applies to the same bad set and generators as HLP3. By (HLP.4)
and injectivity, the complex matrix of \(h_L\) has the same
property. A Jordan chain of length at least two would contradict
it, so this complex matrix is diagonalizable. This can also be
seen without Jordan notation: in each generalized eigenspace,
iteratively apply the stated implication to reduce the least
annihilating power to one.

The minimal polynomial of \(h_L\) is consequently separable over
\(L\). To check descent, the degree of a minimal polynomial is the
least linear dependence among \(1,T,\ldots,T^d\); linear
independence of a finite matrix family is preserved by an injective
field extension, as a nonzero maximal minor remains nonzero.
Thus its scalar extension is the complex minimal polynomial.
A repeated factor over characteristic-zero \(L\) would remain a
repeated root over \(\mathbb C\). Mapping the resulting separable
annihilating polynomial by \(e_c\) gives a separable annihilating
polynomial of \(h_K\). Together with the supplied splitting, this
makes every eligible \(h_K\) semisimple. This is the full original
\texttt{semisimple\allowbreak Tw\_\allowbreak base} under its original no-Jordan input.

The source proves that input by an invariant positive Hermitian
form, its Petersson comparison, and adjoint/commuting double-coset
operators. Those producers are not proved by the foregoing
minimal-polynomial argument and remain explicit when this chapter
is used independently.
% END HLP16

% BEGIN HLP17
\paragraph{HLP17: projection commutes with the algebra and preserves the given lattice.}
\label{claim:HLP17}
Assume \texttt{Semisimple\allowbreak Tw} and \texttt{Hecke\allowbreak Int\allowbreak Tw} on the actual
base, as supplied by HLP16 and HLP8 when their hypotheses hold.
On a simultaneous generalized space \(W_a\), each eligible
semisimple generator acts by its eigenvalue \(\psi_a(h_K)\):
a diagonalizable operator has no larger generalized eigenspaces.
By addition and multiplication, every \(t\in A\) acts on
\(W_a\) as \(\psi_a(t)\). In particular both
\(P_{\mathfrak m}\) and \(Q_{\mathfrak m}\) are \(A\)-stable,
and uniqueness of their direct-sum decomposition gives
\[
                  p_{\mathfrak m}t=tp_{\mathfrak m}.  \tag{HLP.15}
\]
HLP8 gives \(A\)-stability of \(L_0\). Applying (HLP.15) to its
finite sums proves that \(A\) preserves
\(N_{\mathfrak m}=p_{\mathfrak m}L_0\). Restriction therefore
produces an actual \(O\)-algebra homomorphism
\[
          r:A\longrightarrow\operatorname{End}_O(N_{\mathfrak m}).
                                                               \tag{HLP.16}
\]
The equality \(N_{\mathfrak m}=p_{\mathfrak m}L_0\) follows from
(HLP.1) and (HLP.10), so no assertion that a projection preserves
some other unprojected lattice is being used.
% END HLP17

% BEGIN HLP18
\paragraph{HLP18: descent to the exact eigenvalue-function algebra.}
\label{claim:HLP18}
Evaluation on the points of HLP10 gives
\(\operatorname{ev}:A\to O^{X_{\mathfrak m}}\).
Its range is exactly \(T_{\mathfrak m}\): the image contains every
defining function \(h_T(v,k)\), and conversely any element of
\(A\) is a finite polynomial in the generating \(h_K(v,k)\)
with coefficients in \(O\), whose image is the same polynomial
in the \(h_T(v,k)\). Thus it is a surjective map onto this literal
subalgebra.

If \(t\in\ker\operatorname{ev}\), its scalar \(\psi_a(t)\)
vanishes on every \(W_a\) in \(P_{\mathfrak m}\). HLP17 shows
that \(t\) kills \(P_{\mathfrak m}\), hence kills its submodule
\(N_{\mathfrak m}\). Therefore (HLP.16) factors uniquely through
\[
 \operatorname{act}:T_{\mathfrak m}
                       \longrightarrow\operatorname{End}_O(N_{\mathfrak m}).
                                                               \tag{HLP.17}
\]
For every eligible generator its defining identity is exactly
\[
 \operatorname{act}(h_T(v,k))(n)
                    =h_K(v,k)n\quad\text{in }V.       \tag{HLP.18}
\]
This is descent on the given module, with no extension or quotient
of its generators beyond the original \texttt{TOf} construction.
% END HLP18

% BEGIN HLP19
\paragraph{HLP19: the descended action is faithful.}
\label{claim:HLP19}
Suppose \(f\in T_{\mathfrak m}\) has zero action in (HLP.17),
and choose \(t\in A\) mapping to \(f\). It kills
\(N_{\mathfrak m}\). By the full-span conclusion in HLP13 it
kills all of \(P_{\mathfrak m}\). For each
\(\psi_a\in X_{\mathfrak m}\), choose a nonzero vector in
\(W_a\); HLP17 says that its image is \(\psi_a(t)\) times that
vector. Hence \(\psi_a(t)=0\) in \(K\), and injectivity
\(O\hookrightarrow K\) gives \(f(\psi_a)=0\) in \(O\).
Thus the function \(f\) is zero. Applying this to a difference
proves injectivity of \(\operatorname{act}\).

Equations (HLP.17)--(HLP.18) and injectivity are precisely
\texttt{Act\allowbreak Spec}, at every original local index. Choosing this
witness at each index proves the entire \texttt{Act\allowbreak Ex}. In
particular, faithfulness uses all the supported points and the
actual spanning lattice; it is not deduced merely from the
faithful action of a different algebra on \(V\).
% END HLP19

% BEGIN HLP20
\paragraph{HLP20: the source's chosen action satisfies the same operator equation.}
\label{claim:HLP20}
The source may select any witness \(\texttt{act\allowbreak C}\) of
\texttt{Act\allowbreak Ex}. Its \texttt{Act\allowbreak Spec} includes (HLP.18), so on
\(a\otimes n\in K\otimes_ON_{\mathfrak m}\) we have
\[
 \Psi\bigl((1\otimes\operatorname{actC}(h_T(v,k)))(a\otimes n)\bigr)
       =a\,h_K(v,k)n
       =h_K(v,k)\Psi(a\otimes n).                   \tag{HLP.19}
\]
Every tensor is a finite sum of pure tensors, so (HLP.19) extends
to the whole tensor space. This is the full \texttt{Act\allowbreak HK}; it
holds for the source's actual chosen action, not only for one
isomorphic action constructed in HLP18. It also identifies the
base-changed local operators after the transport used in
\texttt{Nd\allowbreak Spec.act\_\allowbreak h\allowbreak K}.
% END HLP20

% BEGIN HLP21
\paragraph{HLP21: locality and nonzero eigenspaces with all inherited hypotheses checked.}
\label{claim:HLP21}
For the exact \(T_{\mathfrak m}\) in HLP10, all its point values
are congruent modulo \(\mathfrak m_O\). Indeed the characters
have the same residual kernel and are \(O\)-algebra maps; their
residual maps identify the quotient with the same residue field
via the fixed scalars. Every polynomial in their generator values
has this property. HLP10 supplies a finite nonempty point set.
Consequently MR24 applies to this very subalgebra of
\(O^{X_{\mathfrak m}}\), proving it local, with nonunits exactly
the functions whose values all belong to \(\mathfrak m_O\).
Concretely, if a value is a unit then all are; the polynomial
\(\prod_x(t-f(x))\) has unit constant term and gives a polynomial
expression for \(f^{-1}\) inside \(O[f]\).

Each point has nonzero generalized space by definition and HLP10.
The finite-dimensional commuting-space argument in MR25 gives
a nonzero common eigenvector for that point. Under HLP17 one may
also take any nonzero vector of its simultaneous space.
HLP13, HLP14 and HLP20 transfer it to the original localized
module with its actual action. Thus every literal local point has
a nonzero eigenspace, as required by
\texttt{hecke\allowbreak Eigenspace\_\allowbreak ne\_\allowbreak bot\_\allowbreak every\_\allowbreak localisation\_\allowbreak point}.
These applications do not import locality or eigenvector
existence as new unproved inputs: MR24--MR25 prove them, and
HLP10--HLP20 verify all their hypotheses on the present objects.
% END HLP21

% BEGIN HLP22
\paragraph{HLP22: the original coefficient map and relation quotient.}
\label{claim:HLP22}
Let \(O_s=D.\mathrm{what}_2.O\), with its specified ring map
\(O_s\to L\). Every image of an element of \(O_s\) in \(L\) is
integral over \(\mathbb Z\) by the supplied row datum. Mapping
that image's monic integer equation by \(e_c\) shows that its image in \(K\) is
integral over \(O\), and hence belongs to \(O\), by the valuation
argument of HLP6. There is therefore one actual ring map
\[
 s_O:O_s\longrightarrow O,\qquad
  \iota_Os_O=e_c\circ(O_s\to L).                    \tag{HLP.20}
\]
It is unique because \(\iota_O:O\hookrightarrow K\) is injective.
This proves the original \texttt{Os\allowbreak Int} and identifies
\texttt{s\allowbreak Of\allowbreak C}.

Let \(G_j\) be the original finite generator index of
\(D.\mathrm{what}_2.\mathrm{MjLat}(j)\), which is distinct from the
ambient \(\mathrm{MtLat}\) set \(G\) in HLP1. Let
\(R\subset O_s^{G_j}\) be its original relation module, and set
\[
 B_j=O^{G_j}/\operatorname{span}_O\{s_O(r):r\in R\}.    \tag{HLP.21}
\]
Images \(b_g\) of the standard basis span this literal quotient,
and \(\sum_g c_gb_g=0\) if and only if the coefficient vector
belongs to the indicated relation span. Under the source's
\texttt{Os\allowbreak Concrete} hypothesis, the image of \(O_s\) is exactly
the integers in a specified number subfield of \(L\); its
injection identifies it with that Dedekind domain. The relations
are saturated because they are the kernel of the given map to a
characteristic-zero vector space. These are precisely the
hypotheses checked and proved sufficient in MR26: split the
finite projective relation quotient over \(O_s\), base change
the resulting idempotent matrix through the actual \(s_O\), and
obtain a finite torsion-free, hence free, \(O\)-module. Thus the
original \texttt{bj\allowbreak C\_\allowbreak free}, span and relation fields all hold.
% END HLP22

% BEGIN HLP23
\paragraph{HLP23: the same coefficient embedding can include all prescribed finite values.}
\label{claim:HLP23}
Fix the source's \(\texttt{Pl\allowbreak Side}\) comparison
\(\iota_w:\overline{\mathbb Q}_\ell\simeq\mathbb C\) and set
\(E_w=\iota_w^{-1}\epsilon:L\to\overline{\mathbb Q}_\ell\).
The field generated by \(E_w(L)\) over \(\mathbb Q_\ell\) is
finite: take a finite \(\mathbb Q\)-basis of the number field
\(L\); these finitely many images generate it as a field, and
all are algebraic over \(\mathbb Q_\ell\). For the prescribed
finite set \(\texttt{big}(w)\subset\mathbb C\), the elements
\(\iota_w^{-1}(z)\) are also algebraic over
\(\mathbb Q_\ell\). Adjoining them gives a finite extension
\(F_2(w)\) containing both required sets. This proves the exact
\texttt{exists\_\allowbreak fin\allowbreak Ext\_\allowbreak big}, with \(\texttt{big}\) given before
the extension.

If the source chooses a larger finite extension \(K_x\) and its
actual integer ring \(O=\mathcal O_{K_x}\), its fraction field
identifies canonically with \(K_x\). Define
\(\operatorname{iq}:\operatorname{Frac}(O)\hookrightarrow
\overline{\mathbb Q}_\ell\) by this identification and the given
inclusion of \(K_x\). It agrees with the original
\(\mathbb Z_\ell\)-algebra map, satisfies
\[
             \iota_w\operatorname{iq}e_c=\epsilon,   \tag{HLP.22}
\]
and has a preimage for every \(z\in\texttt{big}(w)\) after
applying \(\iota_w\). These follow from the inclusions just
constructed and the defining coefficient map of
\texttt{coef\allowbreak In}; they prove \texttt{iq\allowbreak C\_\allowbreak alg},
\texttt{iq\allowbreak C\_\allowbreak ec} and the required \texttt{iq\_\allowbreak big} clause.
% END HLP23

% BEGIN HLP24
\paragraph{HLP24: the unique integral lift of the original twist character.}
\label{claim:HLP24}
Suppose the original \(L^\times\)-valued character \(\psi_0\)
has values whose images under \(e_c\) are units of \(O\), as
asserted by the specified \texttt{Psi\allowbreak Int}. For every group element
choose the unique unit \(\psi(g)\in O^\times\) mapping to
\(e_c(\psi_0(g))\). Multiplicativity and the identity law follow
by applying the injective map to \(K^\times\). Thus this is the
source's \texttt{units\allowbreak Lift} and \texttt{ps\allowbreak N}, with its exact value
formula. Unit integrality itself can be checked by a finite-order
relation when available: if \(a^N=1\) with \(N>0\), both \(a\) and \(a^{-1}\)
are integral over \(O\), so are in \(O\).

The supplied character comparison \(h\psi_0\) identifies
\(\psi_0\), after the actual \(C.\mathrm{embW}\), with
\(C.\psi'\circ D.\mathrm{ntw.art}\), after the prescribed two
coefficient inclusions. All these inclusions are injective.
Consequently
\[
    \psi(g)=1
       \quad\Longleftrightarrow\quad
        (C.\psi'\circ D.\mathrm{ntw.art})(g)=1.       \tag{HLP.23}
\]
For example an element in the kernel of the actual Artin map
belongs to \(\ker\psi\). This is an equality of kernels of the
specified characters, not a comparison with an independently
chosen finite-order character.
% END HLP24

% BEGIN HLP25
\paragraph{HLP25: order, open kernel, anticyclotomic symmetry and exact ramification.}
\label{claim:HLP25}
The original Artin target is elementwise of \(\ell\)-power order.
If its value on \(g\) is killed by \(\ell^a\), (HLP.23) applied
to \(g^{\ell^a}\) gives \(\psi(g)^{\ell^a}=1\).
Its open kernel lies inside \(\ker\psi\), so the latter is open:
a subgroup containing an open subgroup is a union of its open
cosets. The original Artin anticyclotomic identity says that
conjugation carries its values to their inverses. Applying the
same comparison, or (HLP.23) to the corresponding product, proves
the exact \texttt{Anti\allowbreak Hom} of \(\psi\).

Outside \(Q=C.Q(w)\), the actual Artin map kills inertia; hence
\(\psi\) does too. At each \(v\in Q\), retain the original
assertion that \(C.\psi'\circ\mathrm{art}\) is nontrivial on
inertia at every valuation above \(v\). Equation (HLP.23) gives
the same nontrivial element for \(\psi\). This proves
\texttt{Ram\allowbreak Exactly} on the original set \(Q\). If the original
character is trivial on the decomposition groups above \(\ell\),
(HLP.23) proves the corresponding local triviality of \(\psi\).
The character \(\chi_1\) used in \texttt{mk\allowbreak X} is the product of
\(\psi\) and the prescribed transported base twist. At \(v\in Q\)
that base twist is unramified under the printed admissibility
hypotheses; multiplying by it keeps the exhibited nontrivial
inertia value. This proves the required ramification of
\(\chi_1\) whenever these original arithmetic hypotheses are
supplied. Existence of those arithmetic character hypotheses is
not replaced by this kernel calculation.
% END HLP25

% BEGIN HLP26
\paragraph{HLP26: tame generators agree with the original selected generators.}
\label{claim:HLP26}
At a selected prime \(v=\mathrm{plGw}(k,t)\), define
\(\gamma(v)\) to be the original stored \(\mathrm{tg}(k,t)\).
At every other prime choose the tame generator constructed in
LS19 for its original valuation. The two selected rational primes
\(g_0,g_1\) are distinct, so their sets of primes in the number
field are disjoint: if a prime ideal contained both integers,
Bezout's relation for the distinct rational primes would put
\(1\) in that ideal. Within each family the given map
\(\mathrm{plGw}(k,-)\) is injective. Thus every selected prime
has exactly its prescribed index, and the definition returns its
original generator.

LS19 proves membership in inertia and the exact density statement
modulo the source's explicitly defined wild subgroup for every
remaining prime. The stored generators have those same fields
by hypothesis. Hence \(\gamma\) satisfies the original
\texttt{Tame\allowbreak Gen\allowbreak At} everywhere and satisfies
\(\gamma(\mathrm{plGw}(k,t))=\mathrm{tg}(k,t)\) literally.
This is \texttt{gam\allowbreak N\_\allowbreak gen} and \texttt{gam\allowbreak N\_\allowbreak tg}; in particular the
all-prime construction does not disturb either original selected
family, including the prime3 conventions already stored in \(D\).
% END HLP26

% BEGIN HLP27
\paragraph{HLP27: all fields of the original \texttt{Loc\allowbreak Nd} constructor.}
\label{claim:HLP27}
Here is the exact \texttt{mk\allowbreak X} assembly. Retain a base \(B\),
\texttt{Hecke\allowbreak Std0}, \texttt{Fin\allowbreak Dim\allowbreak V}, \texttt{Mt\allowbreak Loc},
\texttt{Act\allowbreak Ex} and \texttt{Act\allowbreak HK} as above. Also retain the source's
specified twist-ring maps, continuity of \(\chi_1\), split and
disjoint exceptional/\(Q\) sets, \(Q\) away from \(\ell\), and
ramification at \(Q\). Set the local indices to HLP10's literal
\(\texttt{Loc\allowbreak T}\); its \(j,w,\ell,\mu\), coefficient ring and
fraction field are the original projections. Use HLP10's finite
matching fibers for \texttt{locs}, its actual \(X_{\mathfrak m}\)
and \(T_{\mathfrak m}\), HLP12's \(N_{\mathfrak m}\), the source's
chosen faithful action from HLP19, and HLP14's
\(\operatorname{pr}_{Mt}\). HLP15 supplies the span and integral
projection clauses, HLP21 supplies locality and all nonzero
eigenspaces, and HLP22 supplies the actual free \(B_j\) and its
generators. The Hecke generator is exactly \(h_T(v,k)\).
Its generated-algebra clause follows from its definition in
HLP10, expanding all six conditions in \(\operatorname{GenAt}\).

Take \(\gamma\) from HLP26. Set
\(\chi(c)=\chi_1\chi_{\rm cyc}^{-(c-1)}\), so
\(\chi(1)=\chi_1\) by exponent zero. Continuity is the supplied
continuity of \(\chi_1\) and of the actual cyclotomic character;
the defining equation and the \(c=1\) equality follow by
substitution. Use exactly the source's finite exceptional set
\(S\), excluding places above \(\ell\), and its original \(Q\).
The supplied splitting/disjointness assertions give their fields.
Use the original value of \(s_0\): at \(\ell=2\) it is \(v_2(D.\mathrm{host}.QS)\), and otherwise
it is \(\max(1,v_\ell(D.\mathrm{host}.QS))\). These are the
actual prime-factorization exponents, not the rational primes
or a multiple of those exponents. The two required exponent
identities follow by the indicated case split.
The ramification field is HLP25 or the identical supplied assertion.
These assignments account for every field of \texttt{mk\allowbreak X}; they
make no automorphic or Galois representation on \(N_{\mathfrak m}\)
beyond the ones actually specified by this constructor.
% END HLP27

% BEGIN HLP28
\paragraph{HLP28: the original twist record over this same local datum.}
\label{claim:HLP28}
For \(a=(j,w,\mathfrak m)\), set the twist index to \(C.\mathrm{twI}(w)\),
the side to \(C.\mathrm{sd}(w)\), the twist-ring homomorphism to the
specified \(\mathrm{tw}_{j,w}\), and the local character to
HLP24's actual \(\psi_{j,w}\). The original comparison of
\(\mathrm{tw}_{j,w}\) with the chosen side embedding is retained
for every element of \(D.\mathrm{LL0}\). The identity for
\(\chi_1\) is its product definition in HLP27. HLP25 supplies the
\(\ell\)-power order, open kernel, anticyclotomic, exact
ramification and above-\(\ell\) triviality fields, whenever their
specified arithmetic inputs hold. These are precisely the fields
of \texttt{mk\allowbreak T}, on the very \(\texttt{mk\allowbreak X}\) just constructed.

The two coefficient comparisons needed later are retained
separately: \texttt{Tw\allowbreak Alg\allowbreak At} is the equality on the actual
\(\ell\)-adic integer scalars, whereas \texttt{Tw\allowbreak Rd\allowbreak At} supplies
one common overfield comparing the twist map and the entire
original map \(L\to K\) with \(C.\mathrm{embW}\). Agreement on
\(D.\mathrm{LL0}\) alone would not supply \texttt{Tw\allowbreak Rd\allowbreak At}.
This constructor neither drops that stronger requirement nor
constructs a new unrelated common overfield in its place.
% END HLP28

% BEGIN HLP29
\paragraph{HLP29: transport to a fixed pair preserves the ideal, projection and point labels.}
\label{claim:HLP29}
For a local index \(a\) with \(j(a)=j\), \(w(a)=w\), the source's
\texttt{tr\allowbreak O} and \texttt{tr\allowbreak K} are the canonical identity maps
after substituting these two equalities. Use them as
\(\phi_O,\phi_K\) in \texttt{Nd\allowbreak Spec}, and use the same substitution
on \(\Psi\) to obtain \texttt{tr\allowbreak Psi}. Replacing the equal indices
in its defining expression makes \texttt{tr\allowbreak Psi} literally HLP14's
map. Its semilinearity, injectivity and range are therefore HLP14;
its projection error is exactly (HLP.13). No change of maximal
ideal occurs: \(\texttt{mx\allowbreak C}\) is the same \(\mathfrak m\) from
the local index. HLP10 proves its maximality, injectivity on the
matching fiber and exhaustion of all eligible maximal ideals.

The point bijection \texttt{tr\allowbreak X} is the same substitution on
\(X_{\mathfrak m}\); it sends a point to its own character and
preserves the equation
\(h_T(v,k)(x)=x(h_K(v,k))\). Transporting (HLP.19) proves the
original \texttt{act\_\allowbreak h\allowbreak K}. The casts of the rational space in
\texttt{Psi\_\allowbreak pr} are the casts induced by these very equalities,
so the vector there is the original \(i_Vu\). This verifies all
transport fields in \texttt{Nonempty\allowbreak Spec}'s first section and
explains why an arbitrary isomorphism of local vector spaces
would not suffice.
% END HLP29

% BEGIN HLP30
\paragraph{HLP30: the universal relation clause and the specified character equation.}
\label{claim:HLP30}
The last \texttt{Nd\allowbreak Spec.bj\_\allowbreak rel} field is quantified over every
candidate coefficient map \(f:O_s\to O\) whose fraction-field
square agrees with the original \(e_c\). By (HLP.20),
\(\iota_Of=\iota_Os_O\); injectivity gives \(f=s_O\) on every
element. Thus its asserted relation span is exactly (HLP.21),
not a larger span depending on a new coefficient choice. HLP22
proves the required if-and-only-if for all these maps, while its
standard-basis images give \texttt{bj} and \texttt{bj\_\allowbreak span}.

For the \texttt{psi\_\allowbreak eq} field, suppose an element \(y\in L\)
has the prescribed image under \(C.\mathrm{embW}(w)\) equal to
the original \(C.\psi'\)-Artin value. Comparing with the supplied
formula for \(\psi_0\) and using injectivity of
\(C.\mathrm{embW}(w)\) gives \(y=\psi_0(g)\) in \(L\).
HLP24's defining equality now gives exactly the requested image
of \(\psi(g)\) in \(K\). This proves the field for every such
\(y\) without changing \(C.\mathrm{embW}\), the Artin map or
\(C.\psi'\). The ring-algebra and common-overfield readings remain
the actual \texttt{Tw\allowbreak Alg\allowbreak At} and \texttt{Tw\allowbreak Rd\allowbreak At} from HLP28.
% END HLP30

% BEGIN HLP31
\paragraph{HLP31: the complete original \texttt{nonempty\_\allowbreak spec}, with its full arithmetic boundary.}
\label{claim:HLP31}
Take exactly the input list of \texttt{nonempty\_\allowbreak spec}. Its
finite algebraic part consists of the specified base \(B_0\),
choice \(C\), \texttt{Hecke\allowbreak Std0}, \texttt{Fin\allowbreak Dim\allowbreak V}, \texttt{Mt\allowbreak Loc},
finite/free \texttt{Mt\allowbreak Sub}, local \texttt{TOf}, \texttt{Act\allowbreak Ex},
nonzero point eigenspaces, \texttt{Act\allowbreak HK}, \texttt{Os\allowbreak Int} and
free \texttt{Bj\allowbreak C}. HLP7--HLP22 provide these on the original
objects when their printed premises hold. Its remaining inputs
are the actual twist-ring maps and characters, their continuity
and comparisons, all splitting, disjointness, order, kernel,
symmetry, ramification and local triviality assertions described
in HLP25--HLP28; the given \(\texttt{Pl\allowbreak Side},M,\texttt{big}\);
the operator identity \(h_{\rm def}\), \(Q\)-over-real-place and
conjugation compatibility, and \(B_0.Q=C.Q\); the specified
\(\operatorname{iq}\) with its scalar, complex-embedding and
\(\texttt{big}\) properties; the actual character equation of
HLP30; and the following local bridge data at every place above
\(\ell\): the given valuation corresponds to the norm at most
one under its embedding, its coefficient field satisfies the
original \texttt{Good\allowbreak Coef}, and its scalar map agrees with the
same \(\operatorname{iq}\). These bridge fields are retained,
not replaced by the abstract existence of a local completion.

Form \(X=\texttt{mk\allowbreak X}\) and \(T=\texttt{mk\allowbreak T}\) as in
HLP27--HLP28. The \texttt{Nd\allowbreak Spec} record is now obtained as
follows. Its \(\operatorname{iq}\), scalar, complex and finite-value
fields are the supplied ones (or HLP23). Its base is
\texttt{rec\allowbreak Base}, which copies \(B_0\) and uses the supplied
\(h_{\rm def}\), \(Q\)-over-real and conjugation assertions.
Its maximal ideals, coefficient equivalences, \(\Psi\), range,
projection and point correspondence are HLP29. Its \texttt{std}
field is the original \texttt{Hecke\allowbreak Std0} under these identity
transports. Its transported \(Q\)-equality follows by the same index
substitution, while \texttt{Qc\_\allowbreak eq} uses the retained hypothesis
\(B_0.Q=C.Q\): the generic \texttt{mk\allowbreak X} copies \(B_0.Q\).
The twist-index and side equalities are reflexive because
\texttt{mk\allowbreak T} uses \(C\) for those fields.
The action equation is HLP20 and HLP29. Its complex coefficient
map is the given \(\texttt{Pl\allowbreak Side}\) composed with
\(\operatorname{iq}\); (HLP.22), or the supplied identical
condition, is its coherence equation. The \(\psi\) and twist
reading fields are HLP30 and the two original comparison inputs.
Its stored selected-prime generators are HLP26. All local bridge,
norm, good-coefficient and \(\operatorname{iq}\)-compatibility
fields are exactly the supplied local bridge witnesses. Finally
its \(B_j\), span and universally quantified relation fields are
HLP22 and HLP30. Every field of \texttt{Nd\allowbreak Spec} has now been
assigned with its original carrier and maps, which proves
\(\operatorname{Nonempty}(\texttt{Nd\allowbreak Spec}(\mathrm{pl},M,C,
\texttt{big},X,T))\). This is the full original constructor,
without promoting any of its arithmetic antecedents to a
conclusion.
% END HLP31

% BEGIN HLP32
\paragraph{HLP32: one coefficient extension for all four original requirements.}
\label{claim:HLP32}
The higher \texttt{nd\allowbreak Spec\_\allowbreak exists} theorem uses four finite
extensions of the same algebraic \(\ell\)-adic field: \(F_1(j,w)\)
for all Hecke eigenvalues, \(F_2(w)\) for the fixed coefficient
embedding and the already prescribed \texttt{big}, \(F_3(w)\)
for all original local bridges above \(\ell\), and \(F_4(w)\)
for the actual twist-ring map with its full common-overfield
reading. MR23 supplies the first under commutation; HLP23 supplies
the second. Retain the original uniform \(F_3,F_4\) supplier
statements, each valid for every larger finite extension with
the same embedding. Then
\[
 K_x(j,w)=(F_1(j,w)F_2(w))(F_3(w)F_4(w))             \tag{HLP.24}
\]
is a finite extension: products of bases span a compositum, and
all four fields are subfields of the same fixed algebraic closure.
It contains each of the four actual subfields. Define
\texttt{coef\allowbreak In} using \(K_x\) and exactly \(E_w\) of HLP23.
The four supplier conclusions therefore hold simultaneously;
none requires enlarging the field after the local objects have
been selected. In particular, the twist map, all bridges,
\(\operatorname{iq}\), the prescribed complex embedding and
all requested finite values have a common coefficient choice.

This proves the coefficient-compositum join in the original
\texttt{nd\allowbreak Spec\_\allowbreak exists} and identifies its next producer cuts
precisely. To apply HLP31 there one still has to supply the
original \texttt{Mt\allowbreak Input} and compatible geometric/choice data;
the complex Hecke standard facts and its no-Jordan theorem;
the actual common-reading supplier used for \(F_4\); the original
local-bridge and \texttt{Good\allowbreak Coef} supplier used for \(F_3\);
and the stated global-reciprocity, admissibility and conductor
ramification inputs. The finite Hecke algebra, projected lattice,
faithful action and \texttt{nonempty\_\allowbreak spec} constructor are no
longer additional unspecified obligations. No \texttt{Aut\allowbreak Spec},
full \texttt{NBuilt}, or final RH/GRH theorem is asserted by this
join alone.
% END HLP32
